% !TEX root = chapter-standalone.tex

\section{Homomorphisms of actions}

\todotext{At the moment unclear if we really want to keep this section; to be considered in the future. Great would be to add a number of interesting examples; this topic related in particular to coordinate changes, bisimulation, etc. (dynamical systems topics). Also, we note that it foreshadows natural transformations.}


\begin{ctdefinition}[Homomorphism of monoid actions]
    \label{def:monoid-action-morphism}
    Let~$\monoidA$ be a \SY{monoid}, and let~$\tup{\setA, \act_\setA}$ and~$\tup{\setB, \act_\setB}$ be two \SY{actions} of~$\monoidA$, where~$\act_\setA \colon \monoidA \fto \Endof\setA$ and~$\act_\setB \colon \monoidA \fto \Endof\setB$.

    A \maindef{homomorphism of monoid actions} (also called an \emph{equivariant map}) from~$\tup{\setA, \act_\setA}$ to~$\tup{\setB, \act_\setB}$ is a function~$\mapa \colon \setA \sto \setB$ such that, for every element~$\monela \setin \monoidAset$,
    \begin{equation}
        \label{eq:equivariance}
        \act_\setA(\monela) \mthen \mapa = \mapa \mthen \act_\setB(\monela).
    \end{equation}
\end{ctdefinition}

In other words, condition \cref{eq:equivariance} says that it does not matter whether we first act by~$\monela$ and then apply~$\mapa$, or first apply~$\mapa$ and then act by~$\monela$: the result is the same.

\todotext{@JL: Add a commutative diagram here illustrating the equivariance condition.}

\begin{example}
    Let~$\monoidA$ be the \SY{monoid}~$\tup{\natnumbers, +, 0}$.
    An \SY{action} of~$\monoidA$ on a set~$\setA$ amounts to a choice of function~$\mapa \colon \setA \sto \setA$ (the action of~$1 \setin \natnumbers$), since the action of any~$n \setin \natnumbers$ is then determined as the~$n$-fold composition of~$\mapa$ with itself.
    A homomorphism between two such actions~$\tup{\setA, \mapa}$ and~$\tup{\setB, \mapb}$ is then a function~$\mapc \colon \setA \sto \setB$ satisfying~$\mapa \mthen \mapc = \mapc \mthen \mapb$, which means that~$\mapc$ intertwines the two dynamics.
\end{example}

    
    
    


